\chapter{An editor that helps}
\label{ch:editor}

\chapterbrief
  {Set up an editor that understands Java, so mistakes surface as you type.}
  {What a language server does for you.}
  {Autocomplete, red squiggles, and jump-to-definition.}

\section{Why this matters more in Java than in C}

Java is verbose, and a good editor absorbs most of that verbosity: it writes the
imports, completes the method names, and underlines the type error before you
compile.

Without one, you will spend a real fraction of Volume One fixing errors an editor
would have shown you as you typed. With one, Java stops feeling wordy.

\section{VS Code, the pragmatic default}

Free, works everywhere, and its remote-SSH support is genuinely excellent.

\begin{enumerate}
  \item Install from \code{code.visualstudio.com}.
  \item Open the Extensions panel (the four-squares icon) and install
        \textbf{Extension Pack for Java} by Microsoft. That single pack pulls in
        the language server, debugger, Maven and test support.
  \item Open your project: \textbf{File $\rightarrow$ Open Folder}
        $\rightarrow$ \code{\textasciitilde/midnight-brewer}.
\end{enumerate}

The first time it opens a Java project it spends a minute indexing. Let it.

\subsection*{What you get}

\begin{center}\small
\begin{tabular}{@{}L{0.30\linewidth}L{0.62\linewidth}@{}}
\toprule
\thead{Feature} & \thead{Why it matters here} \\
\midrule
Red squiggles      & Type errors before you run \code{mvn}. \\
Ctrl-Space         & Completes method names --- and shows what a class offers, which is how you learn an unfamiliar API. \\
F12 / Ctrl-click   & Jump to a definition. Invaluable when reading Pi4J. \\
Ctrl-Shift-O       & Auto-add missing imports. \\
Run / Debug lens   & A clickable \emph{Run} above every \code{main} and every \code{@Test}. \\
Breakpoints        & Stop mid-execution and inspect variables. Debugging the crypto in Volume One will feel very different with these. \\
\bottomrule
\end{tabular}
\end{center}

\begin{concept}{What a ``language server'' is}
The intelligence is not in the editor. It is a separate program --- for Java,
Eclipse JDT --- that parses your project, understands the types, and answers
questions like \emph{``what methods does this variable have?''}

The editor just displays those answers. The same server backs VS Code, Vim and
Emacs, which is why the experience is comparable across all three now.

It also explains a common confusion: the language server has its own view of your
project, so occasionally it disagrees with Maven. When errors look impossible,
\textbf{Ctrl-Shift-P $\rightarrow$ ``Java: Clean Java Language Server
Workspace''} resolves most of it.
\end{concept}

\section{Editing files on the Pi}

Two approaches, and the first is better than it sounds.

\subsection{Edit locally, deploy (recommended)}

Your code lives on your laptop, git-tracked, in a proper editor. The Pi only ever
receives compiled output. That is what Chapter~\ref{ch:deploy} set up, and it is
the right default: the Pi holds no state you would cry about losing.

\subsection{Remote-SSH, for config files}

Sometimes you need to edit something that only exists on the Pi ---
\code{/etc/ssh/sshd\_config}, a systemd unit. Install the \textbf{Remote --- SSH}
extension, then:

\textbf{Ctrl-Shift-P} $\rightarrow$ \emph{Remote-SSH: Connect to Host}
$\rightarrow$ \code{brewer}

It reuses your \code{\textasciitilde/.ssh/config}, so the alias you made in
Chapter~\ref{ch:sshkeys} appears in the list. You now have a full editor running
against the Pi's filesystem.

\begin{trap}{Remote-SSH on a Zero 2 W}
VS Code installs a server component on the Pi, and on 512\,MB it is heavy ---
sometimes heavy enough to trigger the out-of-memory killer, which takes
\code{sshd} with it.

Use it for quick config edits, not as your main workspace. If the Pi becomes
unresponsive after connecting, that is what happened; power-cycle and prefer
\code{nano} over SSH for small edits.
\end{trap}

\section{If you prefer IntelliJ IDEA}

The Community Edition is free and is, for Java specifically, better than VS Code
--- stronger refactoring, better test integration, a nicer debugger.

Open the project by selecting \code{pom.xml} and choosing \emph{Open as
Project}. It has no remote-SSH equivalent, but with the deploy loop you barely
need one.

Either is fine. Do not spend an evening choosing.

\section{Settings worth changing immediately}

\begin{description}[leftmargin=0pt, style=nextline]
  \item[Show whitespace] Makes stray tabs visible. \emph{Editor: Render
        Whitespace $\rightarrow$ boundary}.
  \item[Rulers at 80 and 100] A visual reminder to keep lines readable.
        \code{"editor.rulers": [80, 100]}.
  \item[Format on save] Consistent formatting, zero effort.
  \item[Font size] If you plan to film any of this, \textbf{16--18\,pt
        minimum}. Nobody can read your normal size on a phone.
\end{description}

\begin{checkpoint}{The editor understands your project}
Open \code{ToolchainCheck.java} and:

\begin{enumerate}
  \item Type \code{System.out.} --- a completion list appears.
  \item Delete a semicolon --- a red squiggle appears immediately, before any
        build.
  \item Put it back, then click \emph{Run} above \code{main} --- output appears
        in a panel.
  \item Open \code{ToolchainCheckTest.java} and click \emph{Run Test} --- three
        green ticks.
\end{enumerate}

If completions do not appear, the language server has not finished indexing, or
it failed to import the Maven project. Reload the window and give it a minute.
\end{checkpoint}
